\section{The experiment, rather than a parameter manifold}
A finite observer does not possess the posterior that appears in its analysis. It sees the current report and its retained label. Its next label must serve all histories that have been pooled into the old one. This distinction becomes decisive when a report has a state-dependent likelihood: the conditional law of the forgotten history changes on observing the report, and a linear transport formula is no longer valid.

The problem considered here starts with a preparation, positive causal report instruments, and a declared family of executable future tests. Their conditional evaluations form a predictive quotient. The squared task fixes a norm on these evaluations. No covariance, Fisher matrix, spectrum, or smooth parameter manifold is supplied as a geometric axiom. The relevant measures are the laws actually produced by the preparation and report protocol, including the conditional laws entering each finite-state update.

Our central result is Theorem~\ref{thm:main}. Its first part is an exact identity valid without a stability assumption. For a normalized posterior map $F(y,u)=U(y,u)/\ell(y,u)$, where $U$ and $\ell$ are affine before normalization, introduce
\[
 \Phi(u)=\int\frac{\norm{U(y,u)}^2}{\ell(y,u)}\,\lambda(dy).
\]
The map is convex. Its Hessian, when differentiation is justified, is the report-law average of the squared derivative of the normalized prediction. If $X$ is the full-history predictive state and $\bar X$ its retained conditional mean, the error propagated through a report is exactly the conditional Jensen defect of $\Phi$. The remaining loss divides into quantization within each entering label and disagreement between the row centroids assigned to a common outgoing label.

This identity separates three questions that static quantization alone cannot answer: optimality within a row, compatibility of rows at one update, and existence of a compatible recursion through an entire experiment. The lower bound treats all Borel randomized observers. The upper bound is a genuine finite-state construction. It uses either calibrated common cells, with stability measured by the integrated likelihood curvature, or an explicitly stable posterior recursion on a global cover. Analytical conditional means do not become free runtime registers.

The second part of the theorem gives a sharp statistical resource law. The certificates are stated before choosing an observer: actual one-step acquisition small balls, a global cover, and one of two verifiable stability conditions. Under common positive report support, an internally timed observer must pay for every reachable phase. The resulting excess is
\[
 \calR_{\rm aut}(n,M)-B_n\asymp(M-n)^{-2/d},\qquad M\ge n+1.
\]
Here $B_n$ is the full-history Bayes risk of the same prescribed experiment. This is not an assertion that arbitrary active sensing or optimal stopping has been solved. The kernels may vary with the prescribed phase, but no external phase is given to the autonomous observer. A checkpoint encoder has access to the complete history once and satisfies the different law $\calR_{\rm cp}(n,M)-B_n\asymp M^{-2/d}$.

The two principal verifications are genuinely informative, non-Haar experiments. The classical model is a positive, invertible hidden-state transition observed through a nondegenerate Gaussian likelihood. Its conditional acquisition dimension is derived from the Gaussian log-odds map. The quantum model uses full-dimensional weak effects followed by a positivity-improving channel. Its report law depends on the entering density operator, and its effects do not share an eigenbasis. A positive congruence inverse verifies its acquisition mass; a cone metric verifies stable execution. A periodically improving version permits noncontracting stages. A second quantum instrument with a pure-state branch gives a genuine rank-changing acquired measure and retains each stratum's physical mass.

Finally, Theorem~\ref{thm:joint} shows why errors and memory cannot be combined by notation alone. An independently acquired bit, its actual report erasure, and an unobservable calibration offset are placed in the same terminal task. Remembering that bit through an internally timed experiment changes the available phase budget. The theorem proves matching upper and all-observer lower bounds, not merely an additive error estimate.

Conditional expectation, convex perspectives, projective contraction, and static quantization are standard tools. The contribution claimed here is their instrument-level causal combination, its all-observer row/pooling converse, and the jointly counted finite-state constructions. The relation to sequential vector quantization, finite-state coding, recursive marginal quantization, and information-rate converses is made precise in Section~\ref{sec:comparison}. The earlier relative-report energy and exact-circle results are retained in full as Companion G; the temporal theory remains in separate complete companions. None of those realizations is used to prove the present general theorem backwards.
